\section{Практическая часть: драйвер Windows}
\label{sec:Chapter4} \index{Chapter4}

В данной главе рассматривается архитектура драйверов Windows и способы обращения
пользовательских приложений и драйверов режима ядра к общему кольцевому буферу.

\subsection{Задача драйвера устройства логирования}

В гостевой ОС Windows диагностические сообщения могут быть сгенерированы:
\begin{itemize}
    \item приложениями и сервисами пользовательского режима;
    \item другими драйверами режима ядра.
\end{itemize}

В рамках данной задачи подразумевается, что диагностические сообщения могут отправлять приложения, сервисы и драйверы, которым гостевая ОС выдала доступ к интерфейсу LogDev. Для хоста эти источники всё равно остаются недоверенными, поэтому границы сообщений и состояние буфера проверяются на стороне эмулятора.

Драйвер играет роль единой точки доступа к кольцевому буферу: он предоставляет интерфейс для записи сообщений в буфер, совершая проверки его целостности, и уведомляет эмулятор через MMIO-регистр \texttt{KICK} (раздел~\ref{sec:ch2-logdev}).

\subsection{Схема взаимодействия с драйвером}

Модель Windows разделяет выполнение на пользовательский режим и режим ядра.
Запросы к устройствам оформляются как IRP (I/O Request Packet), которые
проходят стек драйверов, начиная с I/O Manager~\cite{windows_wdk}.

Схема взаимодействия с драйвером:
\begin{enumerate}
    \item PnP Manager обнаруживает устройство через ACPI-таблицы;
    \item функциональный драйвер создаёт Device Object (структура ядра Windows, предоставляющая интерфейс для I/O-взаимодействия с устройством);
    \item пользовательский код открывает символьную ссылку на устройство\\
    (\texttt{\textbackslash{}.\textbackslash{}LogDev});
    и вызывает \texttt{DeviceIoControl};
    \item драйвер обрабатывает IOCTL и обращается к разделяемой памяти.
\end{enumerate}

\subsection{Модель разрабатываемого драйвера}

В рамках данной задачи драйвер был разработан по модели WDM (Windows Driver Model) в силу меньшего количества накладных расходов. Точкой входа в драйвер в данной модели является функция \texttt{DriverEntry}, инициализирующая структуру \texttt{DRIVER\_OBJECT} и регистрирующая обработчики для различных типов запросов. \texttt{DriverUnload} является обработчиком для освобождения ресурсов, выделенных драйвером.

В \texttt{DriverEntry} регистрируются обработчики для различных типов запросов:
\begin{itemize}
    \item \texttt{MajorFunction[IRP\_MJ\_CREATE]} --- открытие устройства (вызывается при \texttt{CreateFile} или \texttt{OpenFile} в пользовательском режиме);
    \item \texttt{MajorFunction[IRP\_MJ\_CLOSE]} --- закрытие (вызывается при \texttt{CloseHandle} в пользовательском режиме);
    \item \texttt{MajorFunction[IRP\_MJ\_DEVICE\_CONTROL]} --- обработка IOCTL (вызывается при \\ \texttt{DeviceIoControl} с определённым кодом в пользовательском режиме);
    \item \texttt{MinorFunction[IRP\_MN\_START\_DEVICE]} --- запрос, который PnP менеджер отправляет драйверу для запуска устройства;
    \item \texttt{AddDevice} --- данная функция вызывается PnP менеджером, когда устройство обнаруживается. Основной задачей данной функции является создание Device Object и прикрепление его к стеку устройств.
\end{itemize}

\subsection{Обнаружение регистров через ACPI}
\label{sec:acpi-logdev}

Для большей универсальности адрес и размер MMIO-регистров передаются через ACPI-таблицы. Далее драйвер из регистров получает физический адрес и размер кольцевого буфера, выделенного в UEFI.

Напомним, что эмулятор добавляет в таблицы ACPI устройство со следующими параметрами~\cite{stratovirt_repo}:
\begin{itemize}
    \item имя ACPI-узла: \texttt{LOGD};
    \item \texttt{\_HID}: \texttt{STVRT\_LOG};
    \item \texttt{\_CRS}: дескриптор \texttt{Memory32Fixed} для MMIO-региона
    устройства LogDev (база и длина согласованы с картой памяти эмулятора).
\end{itemize}

Драйвер в обработчике \texttt{IRP\_MN\_START\_DEVICE} получает структуру \\ \texttt{IO\_STACK\_LOCATION} и извлекает из неё ресурсы устройства (дескриптор \\ \texttt{CM\_RESOURCE\_LIST}). Из него извлекается диапазон MMIO для
регистров (таких как \texttt{BUFPTR}, \texttt{KICK} и др.).
Далее драйвер получает физический адрес и размер кольцевого буфера и отображает физические страницы в виртуальное адресное пространство при помощи функций \texttt{MmMapIoSpace} / \texttt{MmGetPhysicalAddress} (в зависимости от типа памяти --- MMIO-регистры или RAM) и сохраняет указатель на структуру заголовка буфера.

\subsection{Интерфейс для обращения из пользовательского режима}

Для взаимодействия с буфером из приложений и сервисов пользовательского режима в драйвере есть отдельный IOCTL, принимающий готовое сообщение (NUL-терминированная строка).
Обработчик, выставленный в \texttt{MajorFunction[IRP\_MJ\_DEVICE\_CONTROL]}:
\begin{enumerate}
    \item получает входной буфер и его тип при помощи \texttt{IoGetCurrentIrpStackLocation};
    \item проверяет размер (не превышает размер буфера);
    \item вызывает общую функцию записи в MPSC-буфер (по реализации аналогична той, что в UEFI);
    \item при успехе записывает в MMIO \texttt{KICK} магическое значение \texttt{0xDEADBABA} для уведомления потока читателя.
\end{enumerate}

Коды IOCTL задаются при помощи макроса \texttt{CTL\_CODE} с уникальным для устройства значением \texttt{DeviceType}, кодом \texttt{Function}, параметром \texttt{Method} (в случае данного драйвера поддерживается только \texttt{METHOD\_NEITHER}) и флагом \texttt{Access}, определяющим права доступа~\cite{windows_wdk}.

Windows поддерживает 4 способа доступа к пользовательским данным:
\begin{enumerate}
    \item \texttt{METHOD\_BUFFERED} --- при котором ядро создаёт промежуточный буфер, копирует в него пользовательский буфер и передаёт его драйверу;
    \item \texttt{METHOD\_IN\_DIRECT} --- I/O Manager проверяет пользовательские страницы, \\ предотвращая их сброс на диск, и передаёт заполненную структуру MDL, позволяющую драйверу читать данные напрямую;
    \item \texttt{METHOD\_OUT\_DIRECT} --- то же, что \texttt{METHOD\_IN\_DIRECT}, но для записи из драйвера в пользовательский буфер;
    \item \texttt{METHOD\_NEITHER} --- драйвер напрямую получает виртуальные адреса пользовательских буферов.
\end{enumerate}

Для данной задачи был выбран именно последний метод передачи данных, так как он позволяет достичь максимальной производительности. Разработанный драйвер самостоятельно проверяет буферы на валидность при помощи функций \texttt{ProbeForRead} и \texttt{ProbeForWrite}.

\subsection{Интерфейс для обращения из других драйверов}

При получении сообщения от других драйверов режима ядра нет необходимости проверять буфер на валидность, поэтому был разработан второй IOCTL, доступный только из режима ядра, который возвращает указатель на функцию для записи:
\begin{lstlisting}
typedef NTSTATUS (*LOGDEV_WRITE_FN)(
    PVOID Message,
    ULONG MessageSize
);
\end{lstlisting}

Таким образом, другой драйвер может работать с буфером напрямую, минуя дополнительные проверки адреса буфера. Более того, появляется возможность сэкономить на количестве вызовов драйвера.

Проверка того, из какого режима произошёл вызов осуществляется при помощи поля \texttt{RequestorMode} структуры \texttt{IRP}.
